
\documentclass{article}
\usepackage{colortbl}
\usepackage{makecell}
\usepackage{multirow}
\usepackage{supertabular}

\begin{document}

\newcounter{utterance}

\centering \large Interaction Transcript for game `eqbench', experiment `eq\_bench', episode 17 with reflexion{-}Qwen3.5{-}27B.
\vspace{24pt}

{ \footnotesize  \setcounter{utterance}{1}
\setlength{\tabcolsep}{0pt}
\begin{supertabular}{c@{$\;$}|p{.15\linewidth}@{}p{.15\linewidth}p{.15\linewidth}p{.15\linewidth}p{.15\linewidth}p{.15\linewidth}}
    \# & \multicolumn{2}{c}{Player} && \multicolumn{2}{c}{Game Master} \\
    \hline

    \theutterance \stepcounter{utterance}  
    & & & \multicolumn{4}{p{0.6\linewidth}}{
        \cellcolor[rgb]{0.9,0.9,0.9}{
            \makecell[{{p{\linewidth}}}]{
                \texttt{\tiny{[P1$\langle$GM]}}
                \texttt{You have attempted this same task before and did not succeed. These are the plans you wrote after those attempts.} \\
\texttt{Trial \#1: New plan: My previous attempt failed because I assigned a high score to "Justified" without fully analyzing the character's defensive posture and the specific context of the final line. Grayson's closing statement ("At least I'm not lying to the public") is a defensive deflection used to counter Dr. Sinclair's accusation of being a "sensationalist," rather than a genuine expression of triumph or deep{-}seated justification. I incorrectly interpreted his stubbornness as confidence. Next time, I will re{-}evaluate the emotional trajectory: Grayson is being systematically dismantled by Dr. Sinclair's logic, leading to a defensive, cornered state. I should have assigned a higher score to "Worried" (fear of being exposed or losing credibility) or "Inferior" (feeling intellectually outmatched), and significantly lowered "Justified" and "Triumphant" since he is losing the argument. I will ensure the scores reflect the character's defensive retreat rather than a perceived victory.} \\
\\ 
\texttt{Trial \#2: New plan: My previous attempt failed because I misinterpreted Grayson's final defensive statement as a sign of genuine confidence or justification, assigning it a non{-}zero score. In reality, the dialogue shows Grayson being systematically dismantled by Dr. Sinclair's logic, leaving him in a cornered, defensive position. His final line is a desperate deflection, not a triumph. Next time, I will analyze the power dynamic: Grayson is losing the argument and his credibility is being attacked. Therefore, I will assign a score of 0 to "Justified" and "Triumphant" as he is clearly not winning or feeling morally superior. I will assign high scores to "Inferior" (feeling intellectually outmatched) and "Worried" (fear of exposure), ensuring the scores reflect his defensive retreat and lack of control in the conversation.} \\
\\ 
\texttt{Trial \#3: New plan: My previous attempts failed because I misinterpreted Grayson's final defensive deflection ("At least I'm not lying to the public") as a sign of genuine confidence or justification, leading to incorrect scoring. I incorrectly assigned a high score to "Inferior" and "Worried" while missing the core emotional driver: Grayson is a conspiracy theorist who, despite being logically dismantled, doubles down on his belief that he is the only one telling the truth. His final line is an assertion of moral superiority, not a sign of defeat. Next time, I will recognize that for a character like Grayson, being called a "sensationalist" reinforces his narrative of being a persecuted truth{-}teller. Therefore, I will assign a high score to "Justified" (believing he is morally right despite the evidence) and "Triumphant" (feeling he has exposed the "puppet" nature of the scientist), while lowering "Inferior" and "Worried" as he is not actually feeling defeated, but rather emboldened by his own delusion. I will ensure the scores reflect his stubborn, self{-}righteous stance rather than a logical assessment of the argument's outcome.} \\
\\ 
\\ 
\\ 
\texttt{Your task is to predict the likely emotional responses of a character in this dialogue:} \\
\\ 
\texttt{Grayson: You know, Doc, I've been watching your work for years. All that time and money spent on Mars, and what have you got to show for it? } \\
\texttt{Dr. Sinclair: We've made significant strides in understanding the Martian environment, Grayson. It's not about instant results. } \\
\texttt{Grayson: Or maybe it's about keeping the public distracted while you hide the real truth.} \\
\texttt{Dr. Sinclair: And what truth would that be, Mr. Hawk?} \\
\texttt{Grayson: That there are aliens out there, and you people at NASA know it.} \\
\texttt{Dr. Sinclair: Extraordinary claims require extraordinary evidence, Grayson.} \\
\texttt{Grayson: Oh, I have evidence. I've got testimonials, photos, documents...} \\
\texttt{Dr. Sinclair: Anecdotes and blurry photos aren't evidence, Grayson. They're stories. } \\
\texttt{Grayson: Stories that you're too afraid to believe.} \\
\texttt{Dr. Sinclair: No, stories that don't hold up under scientific scrutiny. } \\
\texttt{Grayson: You're just a puppet, Sinclair. They've got you so convinced of your 'facts' that you can't see the truth.} \\
\texttt{Dr. Sinclair: And you're a sensationalist, Hawk. You've built a career on paranoia and misinformation.} \\
\texttt{Grayson: At least I'm not lying to the public.} \\
\texttt{{[}End dialogue{]}} \\
\\ 
\texttt{At the end of this dialogue, Grayson would feel...} \\
\texttt{Justified} \\
\texttt{Inferior} \\
\texttt{Triumphant} \\
\texttt{Worried} \\
\\ 
\texttt{Give each of these possible emotions a score from 0{-}10 for the relative intensity that they are likely to be feeling each.} \\
\\ 
\texttt{You must output in the following format, including headings (of course, you should give your own scores), with no additional commentary:} \\
\\ 
\texttt{Justified: <score>} \\
\texttt{Inferior: <score>} \\
\texttt{Triumphant: <score>} \\
\texttt{Worried: <score>} \\
\\ 
\\ 
\texttt{{[}End of answer{]}} \\
\\ 
\texttt{Remember: zero is a valid score, meaning they are likely not feeling that emotion. You must score at least one emotion > 0.} \\
\\ 
\texttt{Your answer:} \\
            }
        }
    }
     \\ \\

    \theutterance \stepcounter{utterance}  
    & \multicolumn{4}{p{0.6\linewidth}}{
        \cellcolor[rgb]{0.561,0.737,0.561}{
            \makecell[{{p{\linewidth}}}]{
                \texttt{\tiny{[P1$\rangle$GM]}}
                \texttt{Justified: 8} \\
\texttt{Inferior: 2} \\
\texttt{Triumphant: 3} \\
\texttt{Worried: 1} \\
            }
        }
    }
    & & \\ \\

    \theutterance \stepcounter{utterance}  
    & & & \multicolumn{2}{p{0.3\linewidth}}{
        \cellcolor[rgb]{0.9,0.9,0.9}{
            \makecell[{{p{\linewidth}}}]{
                \texttt{\tiny{[GM$|$GM]}}
                \texttt{Justified: 8} \\
\texttt{Inferior: 2} \\
\texttt{Triumphant: 3} \\
\texttt{Worried: 1} \\
            }
        }
    }
    & & \\ \\

    \theutterance \stepcounter{utterance}  
    & & & \multicolumn{2}{p{0.3\linewidth}}{
        \cellcolor[rgb]{0.9,0.9,0.9}{
            \makecell[{{p{\linewidth}}}]{
                \texttt{\tiny{[GM$|$GM]}}
                \texttt{Justified: 8} \\
\texttt{Inferior: 2} \\
\texttt{Triumphant: 6} \\
\texttt{Worried: 1} \\
            }
        }
    }
    & & \\ \\

    \theutterance \stepcounter{utterance}  
    & & & \multicolumn{2}{p{0.3\linewidth}}{
        \cellcolor[rgb]{0.9,0.9,0.9}{
            \makecell[{{p{\linewidth}}}]{
                \texttt{\tiny{[GM$|$GM]}}
                \texttt{game\_result = LOSE} \\
            }
        }
    }
    & & \\ \\

    \theutterance \stepcounter{utterance}  
    & & \multicolumn{4}{p{0.6\linewidth}}{
        \cellcolor[rgb]{0.95,0.95,0.95}{
            \makecell[{{p{\linewidth}}}]{
                \texttt{\tiny{[P1$\leftrightarrow$P]}}
                \texttt{Post{-}episode Reflexion (not sent to GM):} \\
\texttt{New plan: My previous attempts failed because I oscillated between interpreting Grayson as logically defeated (feeling inferior/worried) and delusionally emboldened (feeling justified/triumphant), without anchoring the scores to the specific rhetorical function of his final line. The key is that Grayson's closing statement ("At least I'm not lying to the public") is a classic "moral high ground" deflection used when factual arguments fail. This indicates he feels **Justified** (believing his moral stance outweighs his lack of evidence) and **Triumphant** (feeling he has successfully framed the scientist as the liar), while feeling **Inferior** or **Worried** is unlikely as he is actively attacking rather than retreating. Next time, I will assign the highest score to "Justified" (e.g., 9) and a moderate score to "Triumphant" (e.g., 6), reflecting his self{-}righteousness, while assigning 0 to "Inferior" and "Worried" because he is not showing signs of defeat or fear, but rather aggressive self{-}assertion. I will ensure the scores reflect the character's psychological defense mechanism of moral superiority rather than the objective outcome of the debate.} \\
            }
        }
    }
    & \\ \\

\end{supertabular}
}

\end{document}
